
% 中英文摘要和关键字

\begin{abstract}\label{abstract}

  随着大语言模型逐步从函数级代码生成走向真实仓库级问题求解，代码修复智能体开始具备处理复杂软件工程任务的潜力。
  但在实际执行中，智能体往往需要经历多轮搜索、阅读、编辑与验证，历史上下文持续累积，
  使推理成本快速膨胀，进而影响系统的响应时延、部署成本与工程可用性。
  因此，如何系统理解代码修复智能体的成本来源，并在尽量保持修复能力的前提下降低推理开销，
  成为迈向实际落地必须解决的问题。

  本文围绕这一问题，面向真实仓库级修复场景构建了完整的实验与分析闭环，
  在统一框架下对代码修复智能体的执行过程进行阶段化拆解，并从上下文组织与工作流引导两个维度，
  系统比较了基线方法、文件级RAG、函数级RAG、语义重排序压缩、工作流技能构建、
  技能经验池构建和三阶段组合压缩等方法的成本与效果表现。
  在此基础上，本文建立了以解决率、平均词元消耗和单位成功成本为核心的联合评估框架，
  以刻画不同优化策略在“成本下降”与“能力保持”之间的真实权衡。

  实验结果表明，代码修复智能体的高成本主要并不来自补丁生成本身，而集中于修复位置定位与代码理解过程。
  其中，代码阅读阶段占总词元消耗约一半，失败轨迹则构成主要成本黑洞。
  相比单纯压缩输入内容，更有效的优化方向是提升搜索定位效率、减少无效探索，并改善智能体的决策质量与工作流稳定性。
  在全部方法中，文件级RAG和技能经验池构建表现出更优的综合性价比，能够在控制成本的同时提升整体解决率；
  相比之下，过度细粒度的裁剪或面向通用文本的压缩策略，容易破坏代码任务所依赖的结构信息，收益有限。

  本文的研究表明，代码修复智能体的成本优化应从“压缩更多词元”转向“减少无效过程”，
  即通过更好的上下文选择与更稳定的求解流程，把计算预算集中到真正有助于问题解决的环节上。
  这一结论为高效代码智能体的设计、评估与部署提供了更清晰的分析视角和方法参考。
  
  
    \sysusetup{
      keywords = {大语言模型, 代码修复智能体, 推理成本优化, 上下文压缩, SWE-bench},
    }
  \end{abstract}
  
  \begin{abstract*}
  
  As large language models move from function-level code generation to real
  repository-level issue resolution, code repair agents are becoming increasingly
  capable of handling complex software engineering tasks. Yet this progress also
  exposes a major practical bottleneck: during multi-step execution, agents repeatedly
  search repositories, read code, edit files, and verify patches while continuously
  accumulating interaction history, which drives inference cost upward and limits
  deployability. Understanding where this cost comes from, and how to reduce it
  without substantially harming repair capability, is therefore essential for
  building usable code repair agents.

  This paper presents a systematic study of the cost problem in repository-level 
  code repair. We build a complete experimental pipeline for real repair tasks and
  analyze agent trajectories with a phase-based view of execution. Under a unified
  framework, we compare a baseline method together with several representative
  optimization strategies from two complementary directions: context organization
  and workflow guidance, including file-level retrieval, function-level retrieval,
  semantic re-ranking compression, workflow skill construction, skill memory
  construction, and a three-stage combined compression method. We further establish
  a joint evaluation framework based on resolve rate, average token consumption,
  and cost per resolved instance, so that cost reduction can be assessed together
  with task effectiveness rather than in isolation.

  The results show that the main cost of code repair agents does not lie in patch
  generation itself, but in the process of locating relevant code and understanding
  repository context. Code reading accounts for roughly half of total token usage,
  while failed trajectories form the dominant source of wasted cost. These findings
  suggest that the most effective optimization path is not simply to compress more
  tokens, but to reduce unproductive exploration and improve decision quality.
  Among the evaluated methods, file-level retrieval and skill memory construction
  achieve the best overall trade-off between cost and capability, whereas overly
  aggressive fine-grained trimming or general-purpose text compression tends to
  damage structural information that repository-level repair depends on.

  Overall, this work argues that cost optimization for code repair agents should
  shift from naive input compression toward process-oriented efficiency: selecting
  better context, guiding more stable workflows, and reducing the cost of failure.
  The study provides an experimental foundation and a high-level analytical
  perspective for the design, evaluation, and deployment of more efficient code
  repair agents.
  
  
    \sysusetup{
      keywords* = {Large Language Model, Code Repair Agent, Inference Cost Optimization, Context Compression, SWE-bench},
    }
  \end{abstract*}
  
